\documentclass{thecours}
\begin{document}
    \bigpart{Les deux types de division}
    \subpart{Division euclidienne}\noindent
    La {\bfseries division euclidienne} est un partage équitable avec un reste.\vskip8pt\noindent
    8 pirates souhaitent se partager un trésor de 250 pièces.
    \subpart{Division décimale}\noindent
    La {\bfseries division décimale} est un partage équitable sans reste.\vskip8pt\noindent
    Si on partage une ficelle de 25 cm en 8 morceaux, combien mesure chaque
    morceaux~?
    \bigpart{Lien avec les fractions}
    \definition{quotient_dividende_diviseur}
    \bigpart{Diviseurs et multiples}
    \definition{diviseur_2}{\bfseries Exemple :}
     Le nombre 3 divise-t-il 42 ?\vskip10pt
     \opidiv{42}3 Donc $\dfrac{42}3=14$ et 14 est entier donc 3 divise 42.
     \vskip5pt
     Le nombre 4 divise-t-il 62~?\vskip10pt
     \opidiv{62}4 Donc $\dfrac{62}4=15+\dfrac24$ qui n'est pas entier donc 4
     n'est pas un diviseur de 62.
     \definition{multiple_2}
     {\bfseries Exemple :} 4 est un diviseur de 20 donc 20 est un multiple de 4.
     \proposition{criteres_de_divisibilite}
     {\bfseries Exemple :} On considère le nombre 4\,605.\vskip5pt
     4\,605 se termine par 5, donc c'est un multiple de 5 mais pas un multiple
     de 10 ou de 2.\vskip5pt
     $4+6+0+5=15$ donc 4\,605 est un multiple de 3 mais pas un multiple
     de 9.
\end{document}